\begin{abstract}
Checking agents' cited passages improves majority vote, but an accurate aggregate can still overturn correct answers. We study \emph{verify, then vote}: a fixed plurality default may be overturned only by an answer whose passage an accepted checker labels as supporting it. On OVERTURN's 1,000 primary questions, the pre-specified scored rule raises accuracy from 77.1\% to 85.0\%; on 526 decisive conflicts, from 71.5\% to 86.5\%. It nevertheless misses its pre-specified harm criterion: the upper confidence bound on overturns of correct majorities is 5.5\%, above the 5\% ceiling. \emph{Verified default first}, which keeps a verified default, was developed after this analysis and specified before a second agent pool answered the same questions, partly reusing checker labels. It reaches 85.9\% on that pool's decisive conflicts, with a harm-rate upper bound of 2.2\%. This tests transfer across agent pools, not generalization to unseen questions. Matched-parameter ablations isolate the gate's reduction in harmful overturns, and withholding or replacing cited passages removes most corrections. The formal property identifies an inspectable failure condition: losing a correct default requires a verified wrong answer; it does not establish that such answers are rare. A strong model judge is more accurate and initially cheaper (93.9\% versus 86.5\%; \$0.014 versus \$0.021 per contested primary decision). The benefit of verify-then-vote is an explicit selection policy over reusable evidence: all 1,100 benchmark decisions replay byte-identically, and leave-one-agent-out decisions reuse sealed labels without a model call.

\medskip
\noindent\textbf{Keywords:} multi-agent LLM systems, arbitration, evidence verification, majority vote, replay, provenance, preregistration
\end{abstract}
